\documentclass[a4paper,10pt]{article}

% Настройки русского языка и кодировок
\usepackage[T2A]{fontenc}
\usepackage[utf8]{inputenc}
\usepackage[russian]{babel}

% Математические пакеты
\usepackage{amsmath, amsfonts, amssymb, amsthm}
\usepackage{geometry}
\usepackage{titlesec}
\usepackage{booktabs}

% Настройка полей
\geometry{left=1cm, right=1cm, top=1.5cm, bottom=1.5cm}

% Настройка заголовков
\titleformat{\section}{\large\bfseries}{}{0em}{}
\titleformat{\subsection}{\normalsize\bfseries}{}{0em}{}

\begin{document}

\begin{center}
    \LARGE \textbf{Уравнения Пуассона и Лапласа на сфере (3D)}
\end{center}

\section{Часть 1: Практическая (Алгоритмы и факты)}

\subsection*{Постановка задачи}
Рассматривается краевая задача для уравнения Лапласа ($\Delta u = 0$) или Пуассона ($\Delta u = f$) в трехмерной области $G \subset \mathbb{R}^3$. В сферической системе координат $(r, \theta, \varphi)$, где $\theta \in [0, \pi]$ — полярный угол (коширота), а $\varphi \in[0, 2\pi)$ — азимутальный угол, оператор Лапласа имеет вид:
\[ \Delta u = \frac{1}{r^2}\frac{\partial}{\partial r}\left( r^2 \frac{\partial u}{\partial r} \right) + \frac{1}{r^2} \hat{\Delta}_{\theta,\varphi} u, \quad \text{где } \hat{\Delta}_{\theta,\varphi} = \frac{1}{\sin\theta}\frac{\partial}{\partial \theta}\left( \sin\theta \frac{\partial}{\partial \theta} \right) + \frac{1}{\sin^2\theta}\frac{\partial^2}{\partial \varphi^2} \]
Оператор $\hat{\Delta}_{\theta,\varphi}$ называется оператором \textbf{Лапласа-Бельтрами} на сфере.

Выделяют три основных типа областей и соответствующие им краевые задачи:
\begin{enumerate}
    \item \textbf{В шаре} ($r < R$): Граничное условие ставится на сфере $r=R$. Неявно предполагается ограниченность решения в центре: $|u(0, \theta, \varphi)| < \infty$.
    \item \textbf{Во внешности шара} ($r > R$): Граничное условие на $r=R$. Неявно предполагается условие ограниченности решения на бесконечности: $|u(\infty, \theta, \varphi)| < \infty$.
    \item \textbf{В шаровом слое} ($R_1 < r < R_2$): Задаются два граничных условия: на внутренней сфере $r=R_1$ и на внешней $r=R_2$.
\end{enumerate}

\textbf{Типы краевых условий} (задаются через производную по направлению внешней нормали $\vec{n}$):
\begin{itemize}
    \item \textit{Дирихле (1-й род):} задана сама функция $u|_{\partial G} = \mu(\theta, \varphi)$.
    \item \textit{Неймана (2-й род):} задана производная $u_n|_{\partial G} = \nu(\theta, \varphi)$.
    \item \textit{3-й род:} $(\alpha u + \beta u_n)|_{\partial G} = \mu(\theta, \varphi)$.
\end{itemize}
\textit{Внимание к знакам нормали:} Для шара внешняя нормаль $\partial_n = \partial_r$. Для внешности шара нормаль направлена к центру ($\partial_n = -\partial_r$). Для шарового слоя на $r=R_1$ имеем $\partial_n = -\partial_r$, а на $r=R_2$ имеем $\partial_n = \partial_r$.

\subsection*{Алгоритм решения (Метод разделения переменных)}

\textbf{Шаг 0. Сведение к однородному уравнению} \\
Если $\Delta u = f(r, \theta, \varphi)$, подбирается любое частное решение $w(r, \theta, \varphi)$. Делается замена $u = v + w$, где $\Delta v = 0$ с пересчитанными краевыми условиями.

\textbf{Шаг 1. Разложение краевых условий по сферическим гармоникам} \\
Граничные условия раскладывается в ряд по полному базису действительных сферических функций $Y_n^k(\theta, \varphi)$ (их вид см. ниже). К примеру, для функции $g(\theta, \varphi)$ имеем:
\[ g(\theta, \varphi) = \sum_{n=0}^\infty \sum_{k=-n}^n a_n^k Y_n^k(\theta, \varphi) \]
Коэффициенты вычисляются через скалярное произведение на единичной сфере $S^2$:
\[ a_n^k = \frac{\langle g, Y_n^k \rangle}{\langle Y_n^k, Y_n^k \rangle}, \quad \text{где } \langle f, g \rangle = \int_0^{2\pi} d\varphi \int_0^\pi f(\theta, \varphi) g(\theta, \varphi) \sin\theta \, d\theta \]

\textbf{Упрощенная схема Шага 1} \\
Если краевое условие содержит тригонометрические функции малых степеней, применяется следующий алгоритм:
\begin{itemize}
    \item \textbf{Шаг 1.1.} Представить $g(\theta, \varphi)$ в виде тригонометрического полинома по углу $\varphi$:
    \[ g(\theta, \varphi) = \sum_{k=0}^K L_k(\theta)\cos k\varphi + \sum_{k=1}^K M_k(\theta)\sin k\varphi \]
    \item \textbf{Шаг 1.2.} Разложить получившиеся функции $L_k(\theta)$ и $M_k(\theta)$ в линейные комбинации присоединенных полиномов Лежандра $P_n^k(\cos\theta)$ с \textit{фиксированным} верхним индексом $k$:
    \begin{itemize}
        \item $L_0(\theta)$ выразить через полиномы $P_0^0, P_1^0, P_2^0, P_3^0, \dots$
        \item $L_1(\theta), M_1(\theta)$ выразить через полиномы $P_1^1, P_2^1, P_3^1, \dots$
        \item $L_2(\theta), M_2(\theta)$ выразить через полиномы $P_2^2, P_3^2, \dots$ и так далее.
    \end{itemize}
    \item \textbf{Шаг 1.3.} \textit{Перегруппировать} полученные слагаемые так, чтобы объединить функции с одинаковым \textit{нижним} индексом $n$. В результате краевое условие представится в виде суммы конечного числа гармоник $Y_{n_i}$, где каждая $Y_n$ — это собранная линейная комбинация базисных функций при фиксированном $n$:
    \[ Y_n(\theta, \varphi) = \sum_{k} \left[ c_{n,k} P_n^k(\cos\theta)\cos k\varphi + d_{n,k} P_n^k(\cos\theta)\sin k\varphi \right]. \]
    Итоговое краевое условие примет компактный вид: $g(\theta, \varphi) = Y_{n_1} + Y_{n_2} + \dots + Y_{n_p}$.
\end{itemize}

\textbf{Шаг 2. Выбор радиальной зависимости} \\
Для гармонической функции $\Delta u = 0$ в общем случае решение ищется в виде (обоснование см. в Части 2):
\begin{itemize}
    \item \textbf{Внутри шара ($r < R$):} Отбрасываются расходящиеся при $r \to 0$ слагаемые.
    \[ u(r, \theta, \varphi) = \sum_{n=0}^\infty \sum_{k=-n}^n A_n^k r^n Y_n^k(\theta, \varphi) \]
    \item \textbf{Во внешности шара ($r > R$):} Отбрасываются расходящиеся при $r \to \infty$ слагаемые.
    \[ u(r, \theta, \varphi) = \sum_{n=0}^\infty \sum_{k=-n}^n \frac{B_n^k}{r^{n+1}} Y_n^k(\theta, \varphi) \]
    \item \textbf{В шаровом слое ($R_1 < r < R_2$):} Присутствуют все слагаемые.
    \[ u(r, \theta, \varphi) = \sum_{n=0}^\infty \sum_{k=-n}^n \left( A_n^k r^n + \frac{B_n^k}{r^{n+1}} \right) Y_n^k(\theta, \varphi) \]
\end{itemize}

\textbf{Упрощенная схема Шага 2 (Для конечной суммы гармоник)} \\
Поскольку правая часть свелась к \textit{конечной} сумме конкретных сгруппированных гармоник $Y_{n_i}$, бесконечные ряды писать не нужно. Решение ищется в виде точно такой же конечной суммы:
\[ u(r, \theta, \varphi) = R_{n_1}(r)Y_{n_1}(\theta, \varphi) + R_{n_2}(r)Y_{n_2}(\theta, \varphi) + \dots + R_{n_p}(r)Y_{n_p}(\theta, \varphi), \]
где радиальная часть $R_n(r)$ подбирается индивидуально под каждую гармонику в зависимости от области: 
\begin{itemize}
    \item \textbf{Внутри шара ($r < R$):} $R_n(r) = A_n r^n$
    \item \textbf{Во внешности шара ($r > R$):} $R_n(r) = \frac{B_n}{r^{n+1}}$
    \item \textbf{В шаровом слое ($R_1 < r < R_2$):} $R_n(r) = A_n r^n + \frac{B_n}{r^{n+1}}$
\end{itemize}

\

\textbf{Шаг 3. Подстановка и поиск констант} \\
В общем случае подставить выбранный ряд в граничное условие (с учетом производных для задачи Неймана/3-го рода) и, пользуясь ортогональностью, вычислить коэффициенты.

\

\textbf{Упрощенная схема Шага 3 (Метод неопределенных коэффициентов)} \\
Сконструированную на Шаге 2 конечную сумму (и её производную по $r$, если требуется) подставляем в левую часть краевого условия. 
Например, для обобщенного условия 3-го рода $\left.(\alpha u + \beta u_r)\right|_{r=R} = g(\theta, \varphi)$ получаем тождество:
\[ \sum_{i=1}^p \Big[ \alpha R_{n_i}(R) + \beta R'_{n_i}(R) \Big] Y_{n_i}(\theta, \varphi) \equiv Y_{n_1}(\theta, \varphi) + Y_{n_2}(\theta, \varphi) + \dots + Y_{n_p}(\theta, \varphi). \]
Поскольку функции $Y_n$ (с разными $n$) линейно независимы, мы просто \textit{приравниваем} буквенные коэффициенты при одинаковых гармониках слева и справа. 
Получается набор независимых простых алгебраических уравнений для каждого присутствующего $n_i$:
\[ \alpha R_{n_i}(R) + \beta R'_{n_i}(R) = 1. \]
Для шарового слоя таких условий будет два (на внутренней и внешней границах), что даст системы $2 \times 2$ на константы $A_{n_i}, B_{n_i}$.

\newpage
\subsection*{Таблицы полиномов и сферических функций}

Полиномы Лежандра $P_n(x)$ записываются при помощи формулы Родрига: $P_n(x) = \frac{1}{2^n n!} \frac{d^n}{dx^n} (x^2 - 1)^n$. \\
Присоединенные функции к полиномам Лежандра: $P_n^k(x) = (1-x^2)^{k/2} \frac{d^k}{dx^k} P_n(x)$, где $k \ge 0$.

\vspace{0.3cm}
\begin{center}
\renewcommand{\arraystretch}{1.5}
\begin{tabular}{ll}
\hline
\multicolumn{2}{c}{\textbf{Полиномы Лежандра $P_n(x)$, где $x = \cos\theta$}} \\
\hline
$n=0$ & $P_0(x) = 1$ \\
$n=1$ & $P_1(x) = x$ \\
$n=2$ & $P_2(x) = \frac{1}{2}(3x^2 - 1)$ \\
$n=3$ & $P_3(x) = \frac{1}{2}(5x^3 - 3x)$ \\
\hline
\end{tabular}
\end{center}

\vspace{0.3cm}
\begin{center}
\renewcommand{\arraystretch}{1.5}
\begin{tabular}{ll}
\hline
\multicolumn{2}{c}{\textbf{Присоединенные полиномы Лежандра $P_n^k(x)$, где $x = \cos\theta$}} \\
\hline
$n=0$ & $P_0^0(x) = 1$ \\
\hline
$n=1$ & $P_1^0(x) = x = \cos\theta$ \\
      & $P_1^1(x) = (1-x^2)^{1/2} = \sin\theta$ \\
\hline
$n=2$ & $P_2^0(x) = \frac{1}{2}(3x^2 - 1) = \frac{1}{2}(3\cos^2\theta - 1)$ \\
      & $P_2^1(x) = 3x(1-x^2)^{1/2} = 3\cos\theta\sin\theta$ \\
      & $P_2^2(x) = 3(1-x^2) = 3\sin^2\theta$ \\
\hline
$n=3$ & $P_3^0(x) = \frac{1}{2}(5x^3 - 3x) = \frac{1}{2}(5\cos^3\theta - 3\cos\theta)$ \\
      & $P_3^1(x) = \frac{3}{2}(5x^2-1)(1-x^2)^{1/2} = \frac{3}{2}(5\cos^2\theta-1)\sin\theta$ \\
      & $P_3^2(x) = 15x(1-x^2) = 15\cos\theta\sin^2\theta$ \\
      & $P_3^3(x) = 15(1-x^2)^{3/2} = 15\sin^3\theta$ \\
\hline
\end{tabular}
\end{center}

Сферические функции $Y_n^k(\theta, \varphi)$ образуют полный ортогональный базис на сфере $S^2$. Для $k=0$: $Y_n^0 = P_n^0(\cos\theta)$. Для $k > 0$: $Y_n^k = P_n^k(\cos\theta)\cos(k\varphi)$ и $Y_n^{-k} = P_n^k(\cos\theta)\sin(k\varphi)$.

\vspace{0.3cm}
\begin{center}
\renewcommand{\arraystretch}{1.5}
\begin{tabular}{cllll}
\hline
\textbf{$n$} & \textbf{$k=0$} & \textbf{$k=\pm 1$} & \textbf{$k=\pm 2$} & \textbf{$k=\pm 3$} \\
\hline
$n=0$ & $Y_0^0 = 1$ & & & \\
\hline
$n=1$ & $Y_1^0 = \cos\theta$ & $Y_1^1 = \sin\theta\cos\varphi$ & & \\
      &                      & $Y_1^{-1} = \sin\theta\sin\varphi$ & & \\
\hline
$n=2$ & $Y_2^0 = \frac{1}{2}(3\cos^2\theta - 1)$ & $Y_2^1 = 3\sin\theta\cos\theta\cos\varphi$ & $Y_2^2 = 3\sin^2\theta\cos 2\varphi$ & \\
      &                                          & $Y_2^{-1} = 3\sin\theta\cos\theta\sin\varphi$ & $Y_2^{-2} = 3\sin^2\theta\sin 2\varphi$ & \\
\hline
$n=3$ & $Y_3^0 = \frac{1}{2}(5\cos^3\theta - 3\cos\theta)$ & $Y_3^1 = \frac{3}{2}(5\cos^2\theta-1)\sin\theta\cos\varphi$ & $Y_3^2 = 15\cos\theta\sin^2\theta\cos 2\varphi$ & $Y_3^3 = 15\sin^3\theta\cos 3\varphi$ \\
      &                                                    & $Y_3^{-1} = \frac{3}{2}(5\cos^2\theta-1)\sin\theta\sin\varphi$ & $Y_3^{-2} = 15\cos\theta\sin^2\theta\sin 2\varphi$ & $Y_3^{-3} = 15\sin^3\theta\sin 3\varphi$ \\
\hline
\end{tabular}
\end{center}

\newpage
\section{Часть 2: Теоретическая (Обоснования и доказательства)}

\subsection*{Вывод сферических гармоник и уравнение Лежандра}
\begin{proof}
Будем искать решение уравнения Лапласа $\Delta u = 0$ в виде разделенных переменных $u(r, \theta, \varphi) = R(r)Y(\theta, \varphi)$. Подставим в оператор Лапласа:
\[ \frac{1}{r^2}(r^2 R')'Y + \frac{1}{r^2}R \hat{\Delta}_{\theta,\varphi} Y = 0 \implies \frac{(r^2 R')'}{R} = -\frac{\hat{\Delta}_{\theta,\varphi} Y}{Y}. \]
Левая часть зависит только от $r$, правая — только от углов. Приравняем их константе $\lambda$. Для угловой части получаем спектральную задачу на сфере:
\[ \hat{\Delta}_{\theta,\varphi} Y(\theta, \varphi) + \lambda Y(\theta, \varphi) = 0. \]
Ищем решение в виде $Y(\theta, \varphi) = \Theta(\theta)\Phi(\varphi)$. Подставим в оператор Лапласа-Бельтрами и разделим на $\Theta\Phi$:
\[ \frac{1}{\Theta \sin\theta} \frac{d}{d\theta}\left(\sin\theta \frac{d\Theta}{d\theta}\right) + \frac{1}{\Phi \sin^2\theta} \frac{d^2\Phi}{d\varphi^2} + \lambda = 0. \]
Умножим на $\sin^2\theta$:
\[ \frac{\sin\theta}{\Theta} \frac{d}{d\theta}\left(\sin\theta \frac{d\Theta}{d\theta}\right) + \lambda\sin^2\theta = -\frac{\Phi''}{\Phi} = m^2. \]
Отсюда $\Phi'' + m^2\Phi = 0$. Условие гладкости и $2\pi$-периодичности (аналогично 2D случаю) требует, чтобы $m$ было целым числом ($m=0, 1, 2, \dots$). Решения имеют вид $\Phi_m(\varphi) = A\cos(m\varphi) + B\sin(m\varphi)$.

Для функции $\Theta(\theta)$ получаем уравнение:
\[ \frac{1}{\sin\theta} \frac{d}{d\theta}\left(\sin\theta \frac{d\Theta}{d\theta}\right) + \left( \lambda - \frac{m^2}{\sin^2\theta} \right)\Theta = 0. \]
Сделаем замену переменной $x = \cos\theta$. Тогда $dx = -\sin\theta d\theta$, и $\frac{d}{d\theta} = -\sin\theta \frac{d}{dx} = -\sqrt{1-x^2} \frac{d}{dx}$. Уравнение преобразуется к \textbf{присоединенному уравнению Лежандра}:
\[ \frac{d}{dx}\left( (1-x^2) \frac{d\Theta}{dx} \right) + \left( \lambda - \frac{m^2}{1-x^2} \right)\Theta = 0 \implies (1-x^2) \Theta'' - 2x \Theta' + \left( \lambda - \frac{m^2}{1-x^2} \right)\Theta = 0. \]

Для физичности решения функция $\Theta(x)$ должна быть конечной (ограниченной) на полюсах сферы, то есть при $x \to \pm 1$. Сначала рассмотрим осесимметричный случай $m=0$. Уравнение принимает вид классического уравнения Лежандра:
\[ (1-x^2) \Theta'' - 2x \Theta' + \lambda \Theta = 0. \]
Ищем решение в виде степенного ряда $\Theta(x) = \sum_{s=0}^\infty c_s x^s$. Подставим ряд в уравнение:
\[ \sum_{s=2}^\infty s(s-1)c_s x^{s-2} - \sum_{s=0}^\infty s(s-1)c_s x^s - \sum_{s=0}^\infty 2s c_s x^s + \lambda \sum_{s=0}^\infty c_s x^s = 0. \]
Сдвинем индекс в первой сумме ($s \to s+2$) и приравняем коэффициенты при одинаковых степенях $x^s$ к нулю. Получим рекуррентное соотношение:
\[ (s+2)(s+1)c_{s+2} - \big[s(s-1) + 2s - \lambda\big]c_s = 0 \implies c_{s+2} = \frac{s(s+1) - \lambda}{(s+2)(s+1)} c_s. \]
Если этот степенной ряд будет бесконечным, то при $s \to \infty$ отношение коэффициентов $\frac{c_{s+2}}{c_s} \to 1$. Из теории рядов известно, что это приводит к расходимости функции на границах интервала сходимости $x = \pm 1$ (решение имеет логарифмическую особенность вида $\ln(1 \pm x)$). 
Чтобы решение было гладким и ограниченным на полюсах сферы, ряд обязан \textbf{оборваться}, то есть превратиться в конечный полином. Это произойдет тогда и только тогда, когда числитель дроби обратится в ноль при некотором целом неотрицательном $s = n$. Отсюда мы получаем строгое условие квантования собственных значений:
\[ \lambda = n(n+1), \quad n = 0, 1, 2, \dots \]
Соответствующие полиномиальные решения (с нормировкой $P_n(1)=1$) называются \textbf{полиномами Лежандра} $P_n(x)$.

В общем случае ($m \neq 0$) делается замена $\Theta(x) = (1-x^2)^{m/2} u(x)$. Подстановка в исходное уравнение показывает, что функция $u(x)$ удовлетворяет уравнению:
\[ (1-x^2)u'' - 2(m+1)xu' + \big[ \lambda - m(m+1) \big] u = 0. \]
Аналогичный анализ сходящегося ряда дает для него рекуррентное соотношение $c_{s+2} = \frac{(s+m)(s+m+1) - \lambda}{(s+2)(s+1)} c_s$. 
Ряд обрывается, если $\lambda = (s+m)(s+m+1)$. Обозначая $n = s+m$ (откуда очевидно, что $n \ge m$), мы снова получаем квантование $\lambda = n(n+1)$.
Функции, получаемые при этом, называются \textbf{присоединенными функциями Лежандра} $P_n^m(x) = (1-x^2)^{m/2} \frac{d^m}{dx^m} P_n(x)$.
Возвращаясь к исходным переменным, мы получаем полный классический базис $Y_n^m(\theta, \varphi) = P_n^m(\cos\theta)(A\cos m\varphi + B\sin m\varphi)$.
\end{proof}

\subsection*{Радиальное уравнение Эйлера (3D)}
\begin{proof}
При $\lambda = n(n+1)$ радиальное уравнение имеет вид:
\[ (r^2 R')' - n(n+1)R = 0 \implies r^2 R'' + 2r R' - n(n+1)R = 0. \]
Это уравнение Эйлера. Сделаем замену $r = e^t$, тогда $r \frac{d}{dr} = \frac{d}{dt}$ и $r^2 \frac{d^2}{dr^2} = \frac{d^2}{dt^2} - \frac{d}{dt}$. Подставляем:
\[ (R_{tt} - R_t) + 2R_t - n(n+1)R = 0 \implies R_{tt} + R_t - n(n+1)R = 0. \]
Характеристическое уравнение: $\mu^2 + \mu - n(n+1) = 0$. 
Его корни легко угадываются: $\mu_1 = n$ и $\mu_2 = -(n+1)$. 
Общее решение: $R_n(t) = A_n e^{nt} + B_n e^{-(n+1)t}$.
Делая обратную замену $e^t = r$, получаем фундаментальную радиальную систему для 3D пространства:
\[ R_n(r) = A_n r^n + \frac{B_n}{r^{n+1}}. \]
\end{proof}

\subsection*{Полнота базиса и почему нет «других» решений}
По теореме из спектральной теории, оператор Лапласа-Бельтрами $\hat{\Delta}_{\theta,\varphi}$ является эллиптическим самосопряженным оператором на компактном многообразии (сфере $S^2$). Из этого следует, что его собственные функции $Y_n^k(\theta, \varphi)$ образуют \textbf{полный ортогональный базис} в гильбертовом пространстве квадратично интегрируемых функций $L_2(S^2)$. 

Это строго доказывает, что любая функция на сфере (в том числе краевые условия) может быть представлена в виде сходящегося ряда по гармоникам. Решение краевой задачи, построенное методом Фурье, исчерпывает всё пространство решений. Поиск в принципиально других видах бессмысленен, так как они в итоге спроецируются на этот же базис.

\subsection*{Замечание: Действительные и квантовомеханические гармоники}
В курсах квантовой механики собственные функции оператора квадрата углового момента $\hat{L}^2 = -\hbar^2 \hat{\Delta}_{\theta,\varphi}$ определяются как комплексные сферические гармоники $Y_l^m \propto P_l^m(\cos\theta) e^{im\varphi}$. 
Главное отличие заключается в том, что комплексные $Y_l^m$ одновременно являются также собственными функциями для оператора проекции момента импульса $\hat{L}_z = -i\hbar \frac{\partial}{\partial \varphi}$.

В уравнениях математической физики (где решение рассматривается как действительная функция) применяются \textit{действительные} сферические гармоники, построенные на $\cos(m\varphi)$ и $\sin(m\varphi)$. Они \textbf{остаются собственными функциями для $\hat{L}^2$} (так как $\lambda = l(l+1)$ не зависит от знака $m$), но \textbf{не являются} собственными для $\hat{L}_z$ (поскольку производная косинуса дает синус, и функция не переходит в себя с постоянным множителем).

\end{document}